\section{Conexión con los dos episodios anteriores}

La entrevista con Su Yu (苏煜) en EP139 trata la historia técnica de los Agent; la entrevista con Luo Fuli (罗福莉) en EP138 trata cómo el paradigma Agent transforma los laboratorios de modelos, y la entrevista con Hong Letong (洪乐潼) en EP137 lleva la frontera de la IA hasta la investigación matemática. Leídos en conjunto, los tres episodios forman un mapa continuo que va de los agentes digitales de propósito general y la organización de los laboratorios de modelos a los sistemas de descubrimiento científico.

\begin{importantbox}{El hilo conductor de los tres episodios leídos en conjunto}
EP139 explica por qué los Agent se convirtieron en un nuevo paradigma; EP138 explica cómo los equipos de modelos reorganizan el posentrenamiento y la capacidad de cómputo en función de los Agent; EP137 explica cómo la IA entra en las matemáticas, un campo altamente formalizado y verificable que, sin embargo, depende de la intuición.
\end{importantbox}

\subsection{Resumen de la sección}

AI for Math es una rama científica de alto valor en la era de los Agent. Reúne en un mismo problema el uso de herramientas, la verificación formal, el razonamiento de largo plazo y la colaboración con expertos humanos.
